\documentclass[aspectratio=169,10pt]{beamer}

\usetheme{Boadilla}
\usecolortheme{dolphin}
\setbeamertemplate{navigation symbols}{}
\setbeamertemplate{footline}[frame number]
\setbeamersize{text margin left=7mm,text margin right=7mm}
\usepackage{amsmath,amssymb,mathtools,booktabs}
\usepackage{microtype}
\usepackage{hyperref}

\definecolor{cbmrblue}{RGB}{27,71,112}
\setbeamercolor{structure}{fg=cbmrblue}
\setbeamercolor{frametitle}{fg=cbmrblue}
\setbeamercolor{block title}{bg=cbmrblue!12,fg=cbmrblue}
\setbeamercolor{block body}{bg=cbmrblue!4}
\setbeamercolor{alerted text}{fg=red!75!black}

\newcommand{\R}{\mathbb{R}}
\newcommand{\E}{\operatorname{E}}
\newcommand{\diag}{\operatorname{diag}}
\newcommand{\blkdiag}{\operatorname{blkdiag}}
\newcommand{\Cov}{\operatorname{Cov}}
\newcommand{\Var}{\operatorname{Var}}
\newcommand{\tr}{\operatorname{tr}}
\newcommand{\bbeta}{\boldsymbol{\beta}}
\newcommand{\bgamma}{\boldsymbol{\gamma}}
\newcommand{\bmu}{\boldsymbol{\mu}}
\newcommand{\btheta}{\boldsymbol{\theta}}
\newcommand{\bxi}{\boldsymbol{\xi}}

\title[CBMR math improvements]{The mathematical core of CBMR improvements}
\subtitle{Selected derivations for discussion}
\author{}
\institute{Based on the expanded CBMR derivations}
\date{}

\begin{document}

\begin{frame}[plain]
  \titlepage
\end{frame}

\begin{frame}{Why focus on the mathematics?}
\begin{columns}[T,onlytextwidth]
\column{0.57\textwidth}
\begin{block}{Goal}
Replace large automatic-differentiation calculations with formulas that expose
the likelihood curvature and its structure.
\end{block}
\begin{itemize}
  \item Fit the same specified likelihood; derive its score and information.
  \item Exploit affine predictors, repeated design patterns, and block structure.
  \item Keep statistical assumptions visible when simplifying computation.
\end{itemize}
\column{0.39\textwidth}
\begin{block}{Scale motivating the work}
On the 21-group diagnosis pool, the original Hessian autodiff intermediate was
estimated at \textbf{28.7 GB} for 8,420 coefficients and 21,789 voxels.
\end{block}
\vspace{2mm}
\small The formulas below target this intermediate, not the intrinsic cost of
every possible covariance output.
\end{columns}
\end{frame}

\begin{frame}{One predictor, two effect structures}
For study $i$, voxel $j$, and a fixed offset $o_{ij}$:
\[
\eta_{ij}=\log\mu_{ij}.
\]
\begin{columns}[T,onlytextwidth]
\column{0.49\textwidth}
\begin{block}{Globally constant (GC)}
\[
\eta_{ij}=o_{ij}+B_j^\top\xi_{g(i)}+z_i^\top\gamma.
\]
The moderator effect $\gamma_q$ is constant over voxels.
\end{block}
\column{0.49\textwidth}
\begin{block}{Spatially varying (SV)}
\[
\eta_{ij}=o_{ij}+\sum_{r=1}^R Z_{ir}B_j^\top\beta_r.
\]
The effect of moderator $r$ at voxel $j$ is $B_j^\top\beta_r$.
\end{block}
\end{columns}
\vspace{1mm}
\small For an affine design row $x_{ij}$, the predictor is
$\eta_{ij}=o_{ij}+x_{ij}^\top\btheta$. Fixed dispersion and offsets are
assumed in the regression derivatives that follow.
\end{frame}

\begin{frame}{Closed-form information: scalar derivative to matrix}
Let $\ell_{ij}(\eta_{ij})$ be a cell log-likelihood, with
$s_{ij}=\partial\ell_{ij}/\partial\eta_{ij}$ and
$w_{ij}=-\partial^2\ell_{ij}/\partial\eta_{ij}^2$.
Because $\eta_{ij}$ is affine in $\btheta$:
\[
U=\nabla_{\btheta}\ell=\sum_{i,j}x_{ij}s_{ij},
\qquad
\boxed{H=-\nabla_{\btheta}^2\ell
=\sum_{i,j}w_{ij}x_{ij}x_{ij}^\top.}
\]
\begin{itemize}
  \item No second derivative of the predictor appears; a nonlinear predictor
        would add a score-weighted curvature term.
  \item Poisson: $s=y-\mu$, $w=\mu$.
  \item Independent NB2: $s=(y-\mu)/(1+\alpha\mu)$ and
        $w^{\rm obs}=\mu(1+\alpha y)/(1+\alpha\mu)^2$.
\end{itemize}
\begin{alertblock}{Sign convention}
$H$ is observed information (negative log-likelihood curvature);
$\nabla^2\ell=-H$.
\end{alertblock}
\end{frame}

\begin{frame}{Poisson: exact compression without losing study allocation}
Under a separable exposure, write
$\mu_{ij}=u_{g(i)j}v_i$, where
$u_{gj}=e^{B_j^\top\xi_g}$ and $v_i=e_i e^{z_i^\top\gamma}$.
Define $Y_{gj}=\sum_{i\in\mathcal I_g}Y_{ij}$ and $n_i=\sum_jY_{ij}$.
\[
\ell=\sum_{g,j}Y_{gj}B_j^\top\xi_g
       +\sum_i n_i z_i^\top\gamma
       -\sum_g\left(\sum_j u_{gj}\right)
                \left(\sum_{i\in\mathcal I_g}v_i\right)+C.
\]
\begin{itemize}
  \item This is an exact rearrangement of the full Poisson likelihood.
  \item The compressed statistics include both voxel totals $Y_{gj}$ and
        study totals $n_i$.
  \item Fitting voxel totals alone drops the conditional allocation term
        $\sum_i n_i\log(v_i/T_g)$ and therefore changes the global-covariate
        information.
\end{itemize}
\begin{block}{Takeaway}
Compression is exact because the omitted allocation structure is retained in
the study totals; it is not simply ``sum counts over studies.''
\end{block}
\end{frame}

\begin{frame}{Independent NB: observed and Fisher curvature differ}
For NB2 with mean $\mu=e^\eta$ and dispersion $\alpha$:
\[
\ell_\eta=\frac{y-\mu}{1+\alpha\mu},\qquad
-\ell_{\eta\eta}
=\underbrace{\frac{\mu(1+\alpha y)}{(1+\alpha\mu)^2}}_
                  {w^{\rm obs}}.
\]
Taking expectation at fixed $\mu,\alpha$ gives
\[
\underbrace{\E(w^{\rm obs})}_{w^{\rm F}}
=\frac{\mu}{1+\alpha\mu},
\quad\text{since }\E(Y)=\mu.
\]
\begin{itemize}
  \item Observed weights depend on the realized counts; Fisher weights do not.
  \item Both assemble through $H=X^\top\diag(w)X$, but they answer different
        curvature questions.
  \item As $\alpha\to0$, both weights tend to the Poisson weight $\mu$.
\end{itemize}
\begin{alertblock}{Structural warning}
For NB, separable means do not generally imply separable information weights;
a Poisson Kronecker shortcut is not automatically valid.
\end{alertblock}
\end{frame}

\begin{frame}{Aggregated NB: moment matching is an approximation}
Suppose independent cells have means $\mu_{ij}=u_jv_i$ and variance
$\mu_{ij}+\alpha\mu_{ij}^2$. With
$s_k=\sum_i v_i^k$:
\[
m_j=\E(Y_{\cdot j})=u_js_1,\qquad
\Var(Y_{\cdot j})=u_js_1+\alpha u_j^2s_2.
\]
Matching an NB variance $m_j+m_j^2/\rho$ yields
\[
\rho=\frac{s_1^2}{\alpha s_2},\qquad
A=\frac{s_1}{\alpha s_2},\qquad
p_j=\frac{u_j}{u_j+A}.
\]
\begin{itemize}
  \item This matches the first two moments, not generally the distribution.
  \item The third-cumulant discrepancy is
        $2\alpha^2u_j^3(s_3-s_2^2/s_1)\geq0$.
  \item Equality holds for equal positive study factors; otherwise the
        convolution is generally not exactly NB.
\end{itemize}
\begin{block}{Interpretation}
Derivatives can be exact for this \emph{working likelihood}, while the
likelihood itself remains an approximation to the summed-cell model.
\end{block}
\end{frame}

\begin{frame}{Aggregated NB: the composite Hessian has inner curvature}
The working negative log-likelihood is first differentiated in coordinates
$(\eta,\rho,A)$; then $\rho$ and $A$ are composed with global coefficients
$\gamma$.
For $a,b$ indexing components of $\gamma$:
\[
\begin{aligned}
\Psi_{\gamma_a\gamma_b}={}&
\Psi_{\rho\rho}\rho_a\rho_b
+\Psi_{\rho A}(A_b\rho_a+\rho_bA_a)
+\Psi_{AA}A_aA_b\\
&+\alert{\Psi_\rho\rho_{ab}+\Psi_AA_{ab}}.
\end{aligned}
\]
\begin{itemize}
  \item The last line is the curvature of the inner maps
        $\rho(\gamma),A(\gamma)$.
  \item It does not generally vanish at an observed regression optimum:
        $\Psi_\rho$ and $\Psi_A$ need not separately be zero.
  \item It vanishes in expectation under the correctly specified matched
        distribution, not simply by setting the observed score to zero.
\end{itemize}
\begin{block}{Proof checkpoint}
Differentiate the composite score $\Psi_\rho\rho_a+\Psi_AA_a$ once more;
both product-rule terms produce the highlighted contribution.
\end{block}
\end{frame}

\begin{frame}{Clustered NB: integrate one shared study effect}
Let $\Lambda_i\sim\operatorname{Gamma}(\nu,\text{rate }\nu)$ be shared by all
voxels in study $i$, and
$Y_{ij}\mid\Lambda_i\sim\operatorname{Poisson}(\Lambda_i\mu_{ij})$.
Writing $n_i=\sum_jy_{ij}$ and $t_i=\sum_j\mu_{ij}$, gamma integration gives
\[
\ell_i=\log\Gamma(n_i+\nu)-\log\Gamma(\nu)+\nu\log\nu
-(n_i+\nu)\log(t_i+\nu)
+\sum_j y_{ij}\eta_{ij}+C(y_i).
\]
For $\bmu_i=(\mu_{i1},\ldots,\mu_{iN})^\top$ and
$c_i=(n_i+\nu)/(t_i+\nu)$, the observed information in the log means is
\[
\boxed{W_i^{\rm obs}
=c_i\left\{\diag(\bmu_i)
-\frac{\bmu_i\bmu_i^\top}{t_i+\nu}\right\}.}
\]
\begin{itemize}
  \item Integrating a shared multiplier induces off-diagonal dependence.
  \item The negative rank-one term is essential; independent NB diagonal
        weights describe a different model.
  \item For an affine study design $X_i$, project directly:
        $H_i=X_i^\top W_i^{\rm obs}X_i$.
\end{itemize}
\end{frame}

\begin{frame}{Schur complement: shared effects couple otherwise separate maps}
Suppose the spatial information is block diagonal, but spatial and global
coefficients interact:
\[
H=\begin{pmatrix}D&F\\F^\top&\mathcal E\end{pmatrix},
\qquad T=D^{-1}F,\qquad S=\mathcal E-F^\top T.
\]
Eliminating the spatial coefficients in $Hx=b$ gives
\[
x_\gamma=S^{-1}(b_\gamma-F^\top D^{-1}b_D),\qquad
x_D=D^{-1}b_D-Tx_\gamma.
\]
Consequently,
\[
\boxed{(H^{-1})_{\beta\beta}=D^{-1}+TS^{-1}T^\top,\qquad
(H^{-1})_{\beta\gamma}=-TS^{-1}.}
\]
\begin{itemize}
  \item The spatial correction has rank at most the number $Q$ of global
        coefficients.
  \item Separate spatial information blocks do \emph{not} imply independent
        fitted maps: shared global effects induce cross-group covariance.
  \item Use solves with $D$ and $S$; a dense inverse is not needed for selected
        contrasts.
\end{itemize}
\small $H^{-1}$ is exact matrix algebra when invertible; using it as estimator
covariance is a local/asymptotic approximation with additional assumptions.
\end{frame}

\begin{frame}{Shared spatial patterns: one local calculation, many blocks}
If pattern $h$ has spatial loading $L_h\in\R^R$, its effective map is
$c_h=J_h\bbeta$, with $J_h=L_h^\top\otimes I_P$.
For pattern curvature $K_h=\nabla_{c_h}^2F_h$:
\[
H_{\beta\beta}
=\sum_hJ_h^\top K_hJ_h
=\sum_h(L_hL_h^\top)\otimes K_h.
\]
\begin{itemize}
  \item Compute local pattern curvature once, then scatter it into all
        coefficient blocks (including cross-terms for overlapping loadings).
  \item For independent cells, $B^\top\diag(w)B$ can be formed by row-weighting
        $B$ rather than constructing $\diag(w)$.
  \item This avoids a parameter-by-pattern-by-voxel autodiff intermediate.
\end{itemize}
\begin{alertblock}{Exact boundary}
The information may have exploitable block or pattern structure, but a
requested dense covariance still requires $O(n_\theta^2)$ storage.
\end{alertblock}
\end{frame}

\begin{frame}{What is proved, and what remains a modeling choice}
\begin{columns}[T,onlytextwidth]
\column{0.49\textwidth}
\begin{block}{Evidence in the repository}
\begin{itemize}
  \item SymPy checks exact symbolic identities for selected derivatives and
        matrix formulas.
  \item Lean/Mathlib checks a subset of the derivations.
  \item Independent numerical tests compare derivatives, solves, and fits.
\end{itemize}
\end{block}
\column{0.49\textwidth}
\begin{block}{Keep these distinctions}
\begin{itemize}
  \item Exact differentiation does not make an approximate likelihood exact.
  \item Inverse information is not an exact finite-sample estimator variance.
  \item Mathematical formulas do not imply every model variant is implemented.
  \item Numerical agreement is not formal proof or inferential calibration.
\end{itemize}
\end{block}
\end{columns}
\vfill
\small Primary source: \texttt{docs/cbmr\_expanded\_appendix.tex} and
\texttt{docs/expanded\_covariance.tex}. The current Lean inventory explicitly
does not certify every handwritten intermediate derivation.
\end{frame}

\begin{frame}{Take-home points}
\begin{enumerate}
  \item The affine log-link turns scalar likelihood curvature into a transparent
        weighted design cross-product.
  \item Correct reductions depend on the likelihood: Poisson compression is
        exact with study totals; aggregated NB is moment matched; clustered NB
        retains shared-study dependence.
  \item Block and Schur algebra preserve uncertainty from shared effects while
        enabling structured solves and selected contrasts.
\end{enumerate}
\vfill
\begin{block}{Suggested discussion}
Which distinction is most consequential for the scientific analysis:
GC versus SV effects, the choice among NB constructions, or joint uncertainty
from shared global covariates?
\end{block}
\end{frame}

\end{document}
